% Generated draft: requires mathematical and source review.
\documentclass[11pt]{article}
\usepackage[margin=25mm]{geometry}
\usepackage{fontspec}
\setmainfont{Latin Modern Roman}
\usepackage{amsmath,amssymb,mathtools,unicode-math}
\setmathfont{Latin Modern Math}
\usepackage{enumitem,needspace,xurl,hyperref}
\setlength{\parindent}{0pt}
\setlength{\parskip}{5pt}
\setlength{\emergencystretch}{3em}
\newcommand{\N}{\mathbb{N}}
\newcommand{\Z}{\mathbb{Z}}
\newcommand{\Q}{\mathbb{Q}}
\newcommand{\R}{\mathbb{R}}
\newcommand{\C}{\mathbb{C}}
\newcommand{\F}{\mathbb{F}}
\newcommand{\A}{\mathbb{A}}
\newcommand{\PP}{\mathbb{P}}
\newcommand{\E}{\mathbb{E}}
\newcommand{\eps}{\varepsilon}
\newcommand{\dd}{\,\mathrm{d}}
\begin{document}
\section*{UnsolvedMath ID: 2100709}
\textbf{Source problem code:} catalogue entry
\subsection{Definitions and mathematical statement}
For a circle action $S^1\curvearrowright M$, its fixed-point set is
\[
M^{S^1}=\{x\in M:t\cdot x=x\text{ for every }t\in S^1\}.
\]
A finite number $N=|M^{S^1}|$ requires hypotheses ensuring finitely many isolated fixed points; a general fixed set may have positive-dimensional components. The supplied sentence gives a map $\mu:M\to S^1$, but does not explicitly specify the action whose fixed points are to be counted.

In the generalized moment-map interpretation, a symplectic circle action with infinitesimal generator $X$ can have a circle-valued moment map whose local real lifts satisfy $d\mu=\iota_X\omega$, up to the chosen sign convention. This differs from an ordinary Hamiltonian circle action with a globally real-valued Hamiltonian $H:M\to\mathbb R$. The research summary's Hamiltonian terminology must not silently replace the source's circle-valued target.

\textbf{Question (Á. Pelayo).} Can studying the spectrum of a quantisation of the circle-valued map $\mu:M\to S^1$ give progress in counting the fixed points of the intended action? More concretely, after fixing the action, a quantisation model, and its spectral data, seek spectral information that determines or constrains the fixed-point count. The original request allows progress rather than demanding an exact reconstruction theorem in a specified class. It does not define the quantisation of a circle-valued observable, the spectral parameter or multiplicities, or whether point counts or counts of components are intended. The supplied summary reports no verified result specifically settling this spectral-counting request.
\subsection{Short English statement}
Use the spectrum of a quantised circle-valued moment map to obtain information about fixed-point counts. The action, finiteness assumptions, and circle-valued quantisation model need specification.
\subsection{Sources}
\begin{itemize}
\item \emph{Open Problems, Questions, and Challenges in Finite-Dimensional Integrable Systems} (2018), Problem 7.9, attributed to Á. Pelayo, \url{https://arxiv.org/abs/1804.03737}; source-TeX extraction, accessed 2026-07-29.
\item The supplied summary mentions Atiyah--Bott equivariant localization and the Duistermaat--Heckman formula as context, without publication or theorem identifiers connecting them to this spectral question.
\item Local registry \texttt{data/unsolvedmath/problems.json}, permanent ID $2100709$, source code catalogue entry, \emph{Open Problems in Integrable Systems --- Integrability and Quantisation}. The record labels this open and OPEN-TRIAGE, with NEEDS\_REVIEW source-status and rights qualifications.
\end{itemize}
\end{document}
